\section*{Complete author low-degree division tools}
\paragraph{Theorem 13.} Let $g,h$ be polynomials in $\mathbb F[x]$. If $L(g)=s$ and $\deg(h)=d$, then both $L(g\mathbin{\mathrm{div}}h)$ and $L(g\mathbin{\mathrm{mod}}h)$ have complexity $O(sd)$.
\begin{proof}Suppose $C$ is a circuit of size $s$ which computes $g$. We split every gate $\Phi$ in $C$ into two gates $\Phi_1$ and $\Phi_2$, to make a new circuit $C'$, which computes both $g\mathbin{\mathrm{div}}h$ and $g\mathbin{\mathrm{mod}}h$. If $\Phi$ is computing some polynomial $\phi$ in $C$, then $\Phi_1$ computes the polynomial $\phi\mathbin{\mathrm{mod}}h$ and $\Phi_2$ computes the polynomial $\phi\mathbin{\mathrm{div}}h$.
The proof is inductive and traverses from bottom to the top. The base case is trivial. At some step, say that we are at a gate $\Phi$. The children gate of $\phi$ are computing polynomials $\alpha$ and $\beta$. Let $\alpha=q_1h+r_1$, $\beta=q_2h+r_2$ and $\phi=qh+r$, where the degrees of $r,r_1,r_2$ are smaller than $d$. So in the new circuit $C'$, we have already computed $r_1,q_1,r_2,q_2$. If $\Phi$ is a $\pm$ gate then it is clear that $r=r_1\pm r_2$ and $q=q_1\pm q_2$. If $\Phi$ is a $\times$ gate then we have:
\begin{align*}r&=(r_1r_2)\mathbin{\mathrm{mod}}h,\\q&=q_1q_2h+r_1q_2+r_2q_1+(r_1r_2)\mathbin{\mathrm{div}}h.\end{align*}
We know that $r$ is a polynomial of degree at most $d-1$. Since $\deg(r_1r_2)\le2d-2$, we get that $\deg((r_1r_2)\mathbin{\mathrm{div}}h)\le d-2$. Therefore, we trivially have that: $L(r)=O(d)$ and $L((r_1r_2)\mathbin{\mathrm{div}}h))=O(d)$. Hence we can compute $r,q$ using additional $O(d)$ many gates. Thus, $C'$ has at most $O(sd)$ many gates. Hence $L(g\mathbin{\mathrm{div}}h)=O(sd)$ (same for $g\mathbin{\mathrm{mod}}h$).\end{proof}
\paragraph{Corollary 14.} For $f,g,h\in\mathbb F[x]$, if $f=g/h$ with $L(g)=s$ and $\deg(h)=d$ then $L(f)=O(sd)$.
\paragraph{Remark 15.} The polynomial $f_d:=1+\cdots+x^d=(x^{d+1}-1)/(x-1)$ has $O(\log d)$ size circuit. This can also be shown via a recursive computation argument. Can we expect both $\mathrm{div}$ and $\mathrm{mod}$ to have $\mathrm{poly}(s,\log d)$-size circuits? We show that it is highly unlikely unless factoring is easy, see Theorem 54 in the original paper for details.
\paragraph{Lemma 16 (same as Appendix E Lemma 70).} Suppose $g=\sum_{i\le d_1}g_ix^i$ and $h=x^{d_2}+\sum_{i<d_2}h_ix^i$, in $\mathbb F[x]$. Suppose $g=hq+r$, with $r=\sum_{i<d_2}r_ix^i$ and $q=\sum_{i\le d_1-d_2}q_ix^i$. Then, there is a circuit of size $O(d_1d_2)$, whose inputs are all $h_i,g_i$ and outputs are all $r_i,q_i$.
\paragraph{Complete Appendix E proof (Lemma 70).} We shall denote the desired circuit by $C_{d_1,d_2}$. So we want:
\begin{multline*}C_{d_1,d_2}(g_0,g_1,\ldots,g_{d_1},h_0,h_1,\ldots,h_{d_2})\\=(r_1,r_2,\ldots,r_{d_2-1},q_0,q_1,\ldots,q_{d_1-d_2}).\end{multline*}
If $d_1<d_2$, we know that $q=0$. Hence:
\[C_{d_2-1,d_2}(g_0,g_1,\ldots,g_{d_1-1},h_0,h_1,\ldots,h_{d_1})=(g_1,g_2,\ldots,g_{d_1-1}).\]
If $d_1>d_2$, we perform a long division step:
\begin{align*}g&\leftarrow g-hx^{d_1-d_2}g_{d_1}\\&=\sum_{i\le d_1-d_2-1}g_ix^i+\sum_{i\ge d_1-d_2}^{d_1-1}(g_i-h_{i-(d_1-d_2)}g_{d_1})x^i.\end{align*}
Note that, we can set $q_{d_1-d_2}=g_{d_1}$. Define:
\[\mathbf g\mathrel{\overset{\mathrm{def}}{=}}(g_0,g_1,\ldots,g_{d_1-d_2-1},g_{d_1-d_2}-h_0g_{d_1},\ldots,g_{d_1-1}-h_{d_2-1}g_{d_1}).\]
Then we have:
\begin{multline*}C_{d_1,d_2}(g_0,g_1,\ldots,g_{d_1},h_0,h_1,\ldots,h_{d_2})\\=(C_{d_1-1,d_2}(\mathbf g,h_0,h_1,\ldots,h_{d_2}),g_{d_1}).\tag{8}\end{multline*}
Hence if $S(d_1,d_2)$ is the size of $C_{d_1,d_2}$ then Equation (8) implies that $S(d_1,d_2)=S(d_1-1,d_2)+2d_2$ and $S(d_2-1,d_2)=2d_2-1$. Therefore $S(d_1,d_2)\le2d_1d_2$.\qed
\section*{Complete author monic transformation (Appendix C)}
Given any polynomial $p(\mathbf x)$ in variables $\mathbf x=(x_1,\ldots,x_n)$, there is a standard trick to make it monic in $x_n$ by applying a linear transformation on the variables: for $\alpha=(\alpha_1,\ldots,\alpha_{n-1})\in\mathbb F^{n-1}$, let $\tau_\alpha:x_i\mapsto\alpha_ix_n+x_i$, for $i\in[n-1]$, and $x_n\mapsto x_n$. Note that $\deg(\tau_\alpha(p))\le\deg(p)$ [it may decrease because of non-trivial cancellations]. It is easy to see that $\tau_\alpha$ is an invertible map. We show that $\tau_\alpha(p)$ is monic in $x_n$, for a random transformation $\tau_\alpha$ i.e. when $\alpha\in\mathbb F^{n-1}$ chosen randomly. In fact, we show that this map can simultaneously make polynomials monic given that the field $\mathbb F$ is sufficiently large.
\paragraph{Lemma 68 (Monic Transformation).} Let $p_1(\mathbf x),\ldots,p_m(\mathbf x)$ be $m$-many polynomial of degree $d$. Let $S\subseteq\mathbb F$ be a finite set. For $\alpha\in S^{n-1}$, picked independently and uniformly at random,
\[\Pr\left[\bigwedge_{i=1}^m\tau_\alpha(p_i(\mathbf x))\text{ is monic in }x_n\right]\ge1-\frac{dm}{|S|}.\]
\begin{proof}Consider the terms of degree $d$ of a non-zero polynomial $p\in\mathbb F[\mathbf x]$. Define the set
\[T:=\{\beta=(\beta_1,\ldots,\beta_n):|\beta|_0=\sum_i\beta_i=d,\text{ and }\operatorname{coef}_{\mathbf x^\beta}(p)\ne0\}.\]
We also denote $\beta'=(\beta_1,\ldots,\beta_{n-1})$, the first $n-1$-coordinates of $\beta$, and similarly $\mathbf x'=(x_1,\ldots,x_{n-1})$. Note that $\tau_\alpha(\mathbf x^\beta)=\alpha^{\beta'}x_n^d+(\text{lower terms in }x_n)$. Observe that the homogeneous component of degree $d$ in $\tau_\alpha(p)$ can be written as $a_{d,p}(\mathbf x)=\sum_{\beta\in T}c_\beta\tau_\alpha(\mathbf x^\beta)$, for some constants $c_\beta$. Trivially, $a_{d,p}$ is a nonzero polynomial, and moreover,
\[a_{d,p}(\alpha)=\left(\sum_{\beta\in T}c_\beta\alpha^{\beta'}\right)x_n^d+(\text{lower terms in }x_n).\]
In order to make $\tau_\alpha(p)$ monic in $x_n$, we want $(\sum_{\beta\in T}c_\beta\alpha^{\beta'})\ne0$. So, define, another polynomial $b_{d,p}(\mathbf x')=(\sum_{\beta\in T}c_\beta\mathbf x'^{\beta'})$. It can have degree atmost $d$. As we want each $\tau_\alpha(p)$ monic where $p=p_m(\mathbf x)$, it suffices to find $\alpha$ such that $\prod_{i\in[m]}b_{d,p_i}(\alpha)\ne0$. Note that $\deg(\prod_{i\in[m]}b_{d,p_i}(\mathbf x))\le d\cdot m$. Thus, when we pick $\alpha$ at random, the probability that $\prod_{i\in[m]}b_{d,p_i}(\alpha)=0$, is at most $\le dm/|S|$, from Lemma 64. Hence, the conclusion follows.\end{proof}
